\chapter{벡터공간과 부분공간}

\chapterintro{이 장에서는 선형대수의 기본 언어인 벡터공간(vector space)과 부분공간(subspace)을 엄밀하게 정리하겠습니다. 좌표벡터뿐 아니라 다항식, 함수, 수열처럼 서로 다른 대상이 같은 공리 체계 안에서 통일된다는 점을 확인하시면 이후 장의 기저, 차원, 선형사상 개념을 훨씬 안정적으로 받아들이실 수 있습니다.}

\chaptergoals{벡터공간 공리를 정확히 진술하고 기본 성질을 증명할 수 있다.}{부분공간 판정법을 적용해 다양한 예를 선형 구조로 분류할 수 있다.}{생성(span), 일차독립(linear independence), 직합(direct sum)의 관계를 설명할 수 있다.}

\sectiontemplate{벡터공간의 공리와 기본 성질}{벡터공간은 ``원소의 모양''이 아니라 ``연산 규칙''으로 정의됩니다. 먼저 공리를 정확히 세운 뒤, 공리에서 자동으로 따라오는 기본 성질을 증명해 두면 이후 계산에서 불필요한 혼동을 줄일 수 있습니다.}{\item 벡터공간 공리를 완전한 형태로 진술할 수 있습니다.\item 영벡터와 역벡터의 유일성을 증명할 수 있습니다.\item $0$과 부호에 관한 표준 항등식을 공리로부터 유도할 수 있습니다.}{\item 체(field)의 기본 성질\item 함수와 집합의 기초 표기}

이 절에서는 체 $K$ 위에서 벡터들의 덧셈과 스칼라곱이 어떤 규칙을 만족해야 하는지 먼저 고정하겠습니다. 여기서 한 번 정리한 기호와 논법은 장 전체에서 그대로 사용하겠습니다.

\begin{definition}
체 $K$와 집합 $V$가 주어져 있고, 두 연산
\[
+:V\times V\to V,\qquad \cdot:K\times V\to V
\]
가 다음 공리를 만족한다고 하자. 임의의 $u,v,w\in V$, 임의의 $a,b\in K$에 대하여
\begin{enumerate}[label=(V\arabic*)]
\item $(u+v)+w=u+(v+w)$
\item $u+v=v+u$
\item 어떤 $0_V\in V$가 존재하여 $u+0_V=u$
\item 각 $u\in V$에 대하여 어떤 $-u\in V$가 존재하여 $u+(-u)=0_V$
\item $a(u+v)=au+av$
\item $(a+b)u=au+bu$
\item $(ab)u=a(bu)$
\item $1_Ku=u$
\end{enumerate}
가 성립하면 $V$를 $K$ 위의 \emph{벡터공간}이라 한다.
\end{definition}

\begin{example}
다음은 모두 벡터공간이다.
\begin{enumerate}[label=(\alph*)]
\item $K^n$ (성분별 덧셈, 성분별 스칼라곱)
\item 다항식공간 $K[t]$
\item 구간 $[a,b]$에서 $K$값 연속함수의 집합 $C([a,b],K)$
\end{enumerate}
세 경우 모두 연산 규칙이 다르지 않고 공리 (V1)--(V8)을 그대로 만족한다.
\end{example}

\begin{theorem}[영벡터와 덧셈 역원의 유일성]
벡터공간 $V$에서 덧셈 항등원 $0_V$는 유일하며, 각 $v\in V$의 덧셈 역원 $-v$도 유일하다.
\end{theorem}

\paragraph{증명 전략.}
유일성 증명은 ``둘 있다고 가정한 뒤 서로 같음을 보이는 방식''이 표준적입니다. 항등원은 정의를 양쪽에 한 번씩 적용하고, 역원은 결합법칙과 항등원 성질을 연쇄적으로 사용하면 됩니다.

\begin{proof}
먼저 항등원의 유일성을 보인다. $0_V,0_V'\in V$가 모두 덧셈 항등원이라고 하자. 그러면
\[
0_V=0_V+0_V' \quad(\text{$0_V'$가 항등원})
\]
이고 동시에
\[
0_V+0_V'=0_V' \quad(\text{$0_V$가 항등원})
\]
이므로 $0_V=0_V'$이다.

이제 역원의 유일성을 보인다. $v\in V$를 고정하고 $w,w'\in V$가 모두 $v$의 역원이라고 하자. 즉
\[
v+w=0_V,\qquad v+w'=0_V.
\]
그러면
\[
w=w+0_V=w+(v+w')=(w+v)+w'=(v+w)+w'=0_V+w'=w'
\]
이므로 $w=w'$이다. 따라서 역원도 유일하다.
\end{proof}

\paragraph{의미 해석.}
이 정리는 기호 $0_V$, $-v$를 ``선택''이 아니라 ``정의 가능한 대상''으로 고정해 줍니다. 따라서 이후 계산에서 영벡터나 역벡터를 바꾸어 적을 여지가 없어지고, 모든 증명이 일관된 언어로 진행됩니다.

\begin{theorem}[스칼라 $0$과 부호에 관한 기본 항등식]
임의의 $a\in K$, $v\in V$에 대하여 다음이 성립한다.
\begin{enumerate}[label=(\roman*)]
\item $0_Kv=0_V$
\item $a0_V=0_V$
\item $(-a)v=-(av)$
\item $a(-v)=-(av)$
\item $(-1_K)v=-v$
\end{enumerate}
\end{theorem}

\paragraph{증명 전략.}
핵심은 분배법칙을 두 번 쓰는 것입니다. 먼저 같은 항이 양변에 반복되는 식을 만든 뒤, 덧셈 역원을 더해 소거합니다. 부호 공식은 ``어떤 벡터의 역원임''을 보여 유일성으로 결론내리면 깔끔합니다.

\begin{proof}
(i) 체의 성질로 $0_K+0_K=0_K$이다. 따라서
\[
0_Kv=(0_K+0_K)v=0_Kv+0_Kv.
\]
$0_Kv$의 덧셈 역원을 $y$라 두면
\[
y+0_Kv=y+(0_Kv+0_Kv)=(y+0_Kv)+0_Kv=0_V+0_Kv=0_Kv.
\]
왼쪽 첫 항은 $y+0_Kv=0_V$이므로 $0_V=0_Kv$.

(ii) $0_V+0_V=0_V$이므로
\[
a0_V=a(0_V+0_V)=a0_V+a0_V.
\]
위와 같은 방식으로 양변에 $a0_V$의 역원을 더하면 $a0_V=0_V$.

(iii)
\[
av+(-a)v=(a+(-a))v=0_Kv=0_V
\]
(마지막 등호는 (i)). 즉 $(-a)v$는 $av$의 덧셈 역원이다. 역원의 유일성에 의해 $(-a)v=-(av)$.

(iv)
\[
av+a(-v)=a(v+(-v))=a0_V=0_V
\]
(마지막 등호는 (ii)). 따라서 $a(-v)$는 $av$의 역원이고, 유일성에 의해 $a(-v)=-(av)$.

(v) (iii)에 $a=1_K$를 대입하면
\[
(-1_K)v=-(1_Kv)= -v
\]
(공리 (V8) 사용).
\end{proof}

\paragraph{의미 해석.}
이 항등식들은 계산의 문법을 통일해 줍니다. 이후 절에서 선형결합을 다룰 때 부호 이동, 영항 제거, 동차식 정리를 별도 설명 없이 수행할 수 있는 이유가 바로 여기서 확보됩니다.

\subsection*{주의}
\begin{warning}
연산이 ``정의되어 있다''는 사실만으로 벡터공간이 되지는 않는다. 공리 (V1)--(V8) 중 하나라도 실패하면 벡터공간이 아니다.
\end{warning}
\begin{warning}
같은 집합이라도 스칼라체를 바꾸면 다른 벡터공간이 된다. 예를 들어 $\C$는 $\C$-벡터공간이면서 동시에 $\R$-벡터공간이다.
\end{warning}

\subsection*{자가진단퀴즈}
\begin{enumerate}[label=\arabic*.]
\item 공리 (V5), (V6) 중 어느 것이 스칼라곱의 선형성을 표현하는지 각각 설명하라.
\item $\R^2$의 부분집합 $\{(x,y)\in\R^2:x+y=1\}$가 벡터공간이 아닌 이유를 공리 관점에서 쓰라.
\item $v\in V$에 대하여 $(-1_K)(-v)=v$임을 위 정리를 이용해 증명하라.
\end{enumerate}

\sectiontemplate{부분공간과 부분공간 판정법}{부분공간은 큰 벡터공간 내부의 작은 선형 세계입니다. 실제 문제에서는 공리 8개를 매번 확인하기보다 판정법을 사용해 빠르게 결론을 얻는 경우가 많습니다.}{\item 부분공간의 정의를 벡터공간의 부분구조로 설명할 수 있습니다.\item 선형결합 닫힘 조건으로 부분공간 여부를 판정할 수 있습니다.\item 교집합과 합공간의 보존 성질을 증명할 수 있습니다.}{\item 벡터공간 공리\item 집합의 교집합과 포함관계}

벡터공간의 부분집합이 다시 벡터공간이 되려면, 연산을 제한했을 때 선형 구조가 유지되어야 합니다. 이 절에서는 가장 자주 쓰는 판정 명제를 엄밀하게 증명하겠습니다.

\begin{definition}
벡터공간 $V$의 부분집합 $W\subset V$가 $V$의 연산을 그대로 물려받아 벡터공간을 이루면 $W$를 $V$의 \emph{부분공간}(subspace)이라 한다.
\end{definition}

\begin{theorem}[부분공간 판정법]
$V$를 $K$ 위의 벡터공간이라 하자. 부분집합 $W\subset V$에 대하여 다음이 동치이다.
\begin{enumerate}[label=(\roman*)]
\item $W$는 $V$의 부분공간이다.
\item $W\neq\varnothing$이고, 임의의 $u,v\in W$, $a,b\in K$에 대하여 $au+bv\in W$이다.
\end{enumerate}
\end{theorem}

\paragraph{증명 전략.}
(i)$\Rightarrow$(ii)는 부분공간 정의를 그대로 쓰면 됩니다. 반대 방향에서는 선형결합 닫힘 하나로 $0$, 덧셈, 역원, 스칼라곱을 모두 복원하고, 나머지 공리는 $V$에서 상속된다는 점을 확인합니다.

\begin{proof}
(i)$\Rightarrow$(ii): $W$가 부분공간이면 벡터공간이므로 공집합이 아니다. 또한 덧셈과 스칼라곱에 닫혀 있으므로 $au, bv\in W$이고 다시 $(au)+(bv)=au+bv\in W$.

(ii)$\Rightarrow$(i): 가정에 의해 어떤 $w_0\in W$가 존재한다. 그러면
\[
0_V=0_Kw_0+0_Kw_0\in W
\]
이므로 영벡터가 포함된다. 또한 $u,v\in W$이면
\[
u+v=1_Ku+1_Kv\in W
\]
이므로 덧셈 닫힘이 성립한다. $u\in W$에 대하여
\[
-u=(-1_K)u+0_Ku\in W
\]
이므로 역원 닫힘도 성립한다. 마지막으로 임의의 $c\in K$, $u\in W$에 대하여
\[
cu+0_Ku\in W
\]
이므로 $cu\in W$가 되어 스칼라곱 닫힘이 성립한다.

결합법칙, 교환법칙, 분배법칙, 단위원 법칙 등은 $V$에서 이미 성립하고 $W$는 같은 연산을 제한해 사용하므로 그대로 유지된다. 따라서 $W$는 벡터공간, 즉 $V$의 부분공간이다.
\end{proof}

\paragraph{의미 해석.}
이 정리는 부분공간 판정을 ``선형결합 하나''로 압축해 줍니다. 실제 계산에서는 공리 전체를 다시 쓰지 않고 (ii) 조건만 확인하면 되므로 작업량이 크게 줄어듭니다.

\begin{theorem}[교집합과 합공간의 안정성]
$V$의 부분공간들에 대해 다음이 성립한다.
\begin{enumerate}[label=(\roman*)]
\item 비어 있지 않은 지표집합 $I$와 부분공간족 $\{W_i\}_{i\in I}$에 대해 $\bigcap_{i\in I}W_i$는 부분공간이다.
\item 부분공간 $U,W\le V$에 대해
\[
U+W:=\{u+w:u\in U,\ w\in W\}
\]
는 부분공간이며, $U$와 $W$를 모두 포함하는 가장 작은 부분공간이다.
\end{enumerate}
\end{theorem}

\paragraph{증명 전략.}
(i)는 공통 원소 조건을 직접 확인하는 방식이 가장 간단합니다. (ii)는 선형결합 닫힘으로 부분공간임을 보이고, 포함관계를 이용해 ``가장 작다''는 최소성까지 마무리합니다.

\begin{proof}
(i) 각 $W_i$는 부분공간이므로 $0_V\in W_i$이다. 따라서 $0_V\in\bigcap_{i\in I}W_i$. 이제 $x,y\in\bigcap_{i\in I}W_i$, $a,b\in K$를 잡으면 모든 $i$에 대해 $x,y\in W_i$이므로 $ax+by\in W_i$. 따라서 $ax+by\in\bigcap_{i\in I}W_i$. 부분공간 판정법에 의해 $\bigcap_{i\in I}W_i$는 부분공간이다.

(ii) 먼저 $0_V=0_U+0_W\in U+W$이므로 공집합이 아니다. $x_1=u_1+w_1$, $x_2=u_2+w_2$ ($u_j\in U$, $w_j\in W$)와 $a,b\in K$에 대해
\[
ax_1+bx_2=(au_1+bu_2)+(aw_1+bw_2).
\]
$U,W$가 부분공간이므로 $au_1+bu_2\in U$, $aw_1+bw_2\in W$. 따라서 $ax_1+bx_2\in U+W$. 부분공간 판정법으로 $U+W$는 부분공간이다.

또한 $u=u+0_W\in U+W$, $w=0_U+w\in U+W$이므로 $U,W\subset U+W$. 이제 $M$을 $U,W$를 포함하는 임의의 부분공간이라 하자. $x=u+w\in U+W$이면 $u,w\in M$이므로 $x\in M$. 따라서 $U+W\subset M$. 즉 $U+W$는 최소 부분공간이다.
\end{proof}

\paragraph{의미 해석.}
교집합은 ``공통 조건''을 강화하는 연산, 합공간은 ``허용 벡터''를 확장하는 연산입니다. 두 연산이 부분공간 범주 안에서 닫혀 있다는 사실은 구조를 조합해 문제를 설계할 때 핵심적인 안정성을 제공합니다.

\begin{example}
\[
W=\{(x_1,x_2,x_3,x_4)\in\R^4:x_1-2x_2+x_3=0,\ x_2+x_4=0\}
\]
를 보자. 이 집합은 두 동차 선형방정식의 해집합이며, 각 조건의 해집합은 부분공간이므로 $W$는 그 교집합으로 부분공간이다.
\end{example}

\subsection*{주의}
\begin{warning}
두 부분공간의 합집합 $U\cup W$는 일반적으로 부분공간이 아니다. 부분공간으로 다루려면 $U+W$를 사용해야 한다.
\end{warning}
\begin{warning}
비동차 해집합 $\{x:Ax=b,\ b\neq 0\}$는 보통 부분공간이 아니다. $0$을 포함하지 않는 경우가 대부분이기 때문이다.
\end{warning}

\subsection*{자가진단퀴즈}
\begin{enumerate}[label=\arabic*.]
\item $U,W\le V$일 때 $U\cap W$가 부분공간임을 부분공간 판정법으로 다시 증명하라.
\item $\R^2$에서 두 직선의 합집합이 부분공간이 되지 않는 예를 하나 구성하라.
\item $A\in M_{m,n}(K)$에 대해 $\{x\in K^n:Ax=0\}$가 부분공간임을 설명하라.
\end{enumerate}

\sectiontemplate{생성공간과 일차독립}{벡터들의 모임이 공간을 얼마나 ``만들어 내는지''를 측정하는 개념이 생성이고, 서로 중복 없이 정보를 담는지를 측정하는 개념이 일차독립입니다. 두 개념을 분리해서 이해하면 기저 개념으로 자연스럽게 이어집니다.}{\item 생성공간의 정의와 최소성 성질을 증명할 수 있습니다.\item 일차독립과 일차종속을 계수 관계로 판정할 수 있습니다.\item 종속 벡터 제거 원리를 설명할 수 있습니다.}{\item 선형결합(linear combination)\item 부분공간 판정법}

생성과 독립은 서로 보완적인 기준입니다. 생성은 ``충분성'', 독립은 ``중복 없음''을 뜻합니다. 이 절에서는 생성공간의 최소성 정리와 종속 판정의 표준 기준을 확보하겠습니다.

\begin{definition}
$S\subset V$에 대하여 $S$의 \emph{생성공간}(span)을
\[
\operatorname{span}(S)
=
\left\{
\sum_{i=1}^{m}a_is_i
\;\middle|\;
m\ge 1,\ a_i\in K,\ s_i\in S
\right\}
\cup\{0_V\}
\]
로 정의한다.
\end{definition}

\begin{definition}
집합 $S\subset V$가 \emph{일차독립}(linearly independent)이라는 것은, 서로 다른 유한개 벡터 $v_1,\dots,v_m\in S$와 스칼라 $a_1,\dots,a_m\in K$가
\[
a_1v_1+\cdots+a_mv_m=0_V
\]
를 만족하면 반드시 $a_1=\cdots=a_m=0_K$가 되는 것을 말한다. 이 조건이 실패하면 일차종속(linearly dependent)이라 한다.
\end{definition}

\begin{theorem}[생성공간의 부분공간성 및 최소성]
$S\subset V$에 대하여 다음이 성립한다.
\begin{enumerate}[label=(\roman*)]
\item $\operatorname{span}(S)$는 $V$의 부분공간이다.
\item $S\subset W\le V$이면 $\operatorname{span}(S)\subset W$이다.
\end{enumerate}
특히 $\operatorname{span}(S)$는 $S$를 포함하는 가장 작은 부분공간이다.
\end{theorem}

\paragraph{증명 전략.}
(i)는 선형결합의 집합이 다시 선형결합에 닫혀 있음을 직접 보이면 됩니다. (ii)는 ``부분공간은 선형결합에 닫혀 있다''는 사실을 한 번에 적용하면 즉시 얻어집니다.

\begin{proof}
(i) 먼저 $0_V\in\operatorname{span}(S)$는 정의에서 자명하다. 이제 $x,y\in\operatorname{span}(S)$, $a,b\in K$를 잡자.
\[
x=\sum_{i=1}^{m}\alpha_is_i,\qquad
y=\sum_{j=1}^{n}\beta_jt_j
\]
($s_i,t_j\in S$)로 쓸 수 있다. 그러면
\[
ax+by
=
\sum_{i=1}^{m}(a\alpha_i)s_i
+
\sum_{j=1}^{n}(b\beta_j)t_j
\]
이므로 $S$의 유한 선형결합이며 따라서 $\operatorname{span}(S)$에 속한다. 부분공간 판정법에 의해 $\operatorname{span}(S)$는 부분공간이다.

(ii) $S\subset W$이고 $W\le V$라 하자. $x\in\operatorname{span}(S)$이면 $x=\sum_{i=1}^{m}a_is_i$ ($s_i\in S$)로 쓸 수 있다. 각 $s_i\in W$이고 $W$는 선형결합에 닫혀 있으므로 $x\in W$. 따라서 $\operatorname{span}(S)\subset W$.
최소성은 (i), (ii)를 합치면 즉시 따른다.
\end{proof}

\paragraph{의미 해석.}
생성공간은 ``주어진 벡터들만으로 도달 가능한 전체 영역''을 정확히 표현합니다. 최소성까지 함께 확보되므로, 새로운 부분공간을 정의할 때 사실상 표준 형식으로 사용할 수 있습니다.

\begin{theorem}[유한 벡터열의 일차종속 판정]
유한 벡터열 $(v_1,\dots,v_n)$에 대하여 다음이 동치이다.
\begin{enumerate}[label=(\roman*)]
\item $(v_1,\dots,v_n)$은 일차종속이다.
\item 어떤 $j$가 존재하여 $v_j\in\operatorname{span}(\{v_i:i\neq j\})$이다.
\end{enumerate}
\end{theorem}

\paragraph{증명 전략.}
(i)$\Rightarrow$(ii)는 비자명한 선형관계에서 계수가 0이 아닌 항 하나를 골라 푸는 방식입니다. (ii)$\Rightarrow$(i)는 해당 식을 한쪽으로 이항해 즉시 비자명한 선형관계를 구성하면 됩니다.

\begin{proof}
(i)$\Rightarrow$(ii): 일차종속이므로
\[
a_1v_1+\cdots+a_nv_n=0_V
\]
를 만족하는 계수 $a_1,\dots,a_n$가 존재하고, 모두 $0$은 아니다. 따라서 어떤 $j$에 대해 $a_j\neq 0$. 그러면
\[
v_j
=
-\sum_{i\neq j}a_j^{-1}a_iv_i
\]
가 되어 $v_j\in\operatorname{span}(\{v_i:i\neq j\})$.

(ii)$\Rightarrow$(i): 어떤 $j$에 대해
\[
v_j=\sum_{i\neq j}c_iv_i
\]
라 하자. 이를 이항하면
\[
\sum_{i\neq j}(-c_i)v_i+1\cdot v_j=0_V.
\]
$j$번째 계수 $1$이 0이 아니므로 비자명한 선형관계가 존재한다. 따라서 일차종속이다.
\end{proof}

\paragraph{의미 해석.}
이 정리는 ``종속이면 불필요한 벡터가 있다''는 직관을 엄밀하게 만듭니다. 이후 기저를 구성할 때 종속 벡터를 제거하는 알고리즘의 논리적 근거가 됩니다.

\begin{example}
$\R^3$에서
\[
v_1=(1,0,1),\quad v_2=(0,1,1),\quad v_3=(1,1,2)
\]
를 두면 $v_3=v_1+v_2$이므로 $(v_1,v_2,v_3)$는 일차종속이다. 따라서 $\operatorname{span}\{v_1,v_2,v_3\}=\operatorname{span}\{v_1,v_2\}$이다.
\end{example}

\subsection*{주의}
\begin{warning}
$\operatorname{span}(S)=S$가 항상 성립하는 것은 아니다. 일반적으로 $S$는 생성공간보다 작고, 선형결합을 통해 원소가 새로 생긴다.
\end{warning}
\begin{warning}
일차독립 판정에서는 반드시 동차식 $\sum a_iv_i=0$을 사용해야 한다. 비동차식으로 판단하면 결론이 왜곡된다.
\end{warning}

\subsection*{자가진단퀴즈}
\begin{enumerate}[label=\arabic*.]
\item $S=\{(1,0),(1,1)\}\subset\R^2$에 대해 $\operatorname{span}(S)$를 구하라.
\item 벡터열 $(1,2,3),(2,4,6),(0,1,1)$의 일차종속 여부를 판정하라.
\item ``$v\in\operatorname{span}(S)$이고 $S$가 일차독립이면 $S\cup\{v\}$는 일차종속이다''를 증명하라.
\end{enumerate}

\sectiontemplate{부분공간의 합과 직합}{부분공간을 합쳐 더 큰 공간을 만드는 과정은 구조 해석의 핵심 도구입니다. 특히 직합은 분해가 유일한 경우를 가리키며, 좌표화와 사영 개념으로 이어지는 중요한 관문입니다.}{\item 합공간 $U+W$의 의미를 생성 관점에서 해석할 수 있습니다.\item 직합의 유일분해 조건을 정확히 증명할 수 있습니다.\item 예제를 통해 직합 여부를 판정할 수 있습니다.}{\item 부분공간 판정법\item 생성공간과 일차독립}

이 절에서는 부분공간 두 개를 결합하는 두 가지 관점, 즉 ``그냥 합''과 ``유일 분해가 가능한 합''을 구분하겠습니다. 이 구분은 다음 장의 차원 공식과 선형사상 분해에서 반복적으로 사용됩니다.

\begin{definition}
부분공간 $U,W\le V$에 대하여
\[
U+W=\{u+w:u\in U,\ w\in W\}
\]
를 \emph{합공간}이라 한다. 또한 $V=U+W$이면서 $U\cap W=\{0_V\}$이면
\[
V=U\oplus W
\]
라 쓰고 \emph{직합}(direct sum)이라 한다.
\end{definition}

\begin{theorem}[직합의 유일분해 판정]
$V=U+W$라 하자. 다음이 동치이다.
\begin{enumerate}[label=(\roman*)]
\item 임의의 $v\in V$가 $v=u+w$ ($u\in U,\ w\in W$) 꼴로 유일하게 표현된다.
\item $U\cap W=\{0_V\}$.
\end{enumerate}
\end{theorem}

\paragraph{증명 전략.}
(i)$\Rightarrow$(ii)는 영벡터의 두 표현을 비교하면 됩니다. (ii)$\Rightarrow$(i)는 두 표현의 차를 취해 교집합 원소가 됨을 보이고, 교집합이 자명하다는 가정으로 동일성을 강제합니다.

\begin{proof}
(i)$\Rightarrow$(ii): $z\in U\cap W$를 잡자. 그러면
\[
0_V=0_U+0_W=z+(-z)
\]
이고 $z\in U$, $-z\in W$이므로 $0_V$의 두 분해가 된다. 유일성에 의해 $z=0_U=0_V$. 따라서 $U\cap W=\{0_V\}$.

(ii)$\Rightarrow$(i): $V=U+W$이므로 존재성은 자명하다. 유일성을 보이자. $v=u_1+w_1=u_2+w_2$라 하자($u_i\in U,\ w_i\in W$). 그러면
\[
u_1-u_2=w_2-w_1.
\]
왼쪽은 $U$에, 오른쪽은 $W$에 속하므로 공통 원소이며
\[
u_1-u_2\in U\cap W=\{0_V\}.
\]
따라서 $u_1=u_2$, 이어서 $w_1=w_2$. 유일성이 성립한다.
\end{proof}

\paragraph{의미 해석.}
직합은 단순한 합이 아니라 ``중복 없는 결합''입니다. 한 벡터를 두 성분으로 나누는 방식이 하나로 고정되므로, 복잡한 대상을 독립한 두 부분으로 분해해 해석할 수 있습니다.

\begin{theorem}[합공간의 생성 표현]
임의의 부분공간 $U,W\le V$에 대하여
\[
U+W=\operatorname{span}(U\cup W)
\]
가 성립한다.
\end{theorem}

\paragraph{증명 전략.}
한쪽 포함은 $u+w$를 $U\cup W$의 선형결합으로 바로 쓰면 됩니다. 반대 포함은 임의의 선형결합을 ``$U$쪽 항의 합 + $W$쪽 항의 합''으로 분리하면 됩니다.

\begin{proof}
먼저 $x\in U+W$라 하자. 그러면 $x=u+w$ ($u\in U$, $w\in W$). $u,w\in U\cup W$이므로
\[
x=1_K\cdot u+1_K\cdot w\in\operatorname{span}(U\cup W).
\]
따라서 $U+W\subset\operatorname{span}(U\cup W)$.

반대로 $y\in\operatorname{span}(U\cup W)$라 하자. 그러면
\[
y=\sum_{i=1}^m a_i z_i,\qquad z_i\in U\cup W.
\]
지표를 $I_U=\{i:z_i\in U\}$, $I_W=\{i:z_i\in W\}$로 나누면
\[
y=\sum_{i\in I_U}a_iz_i+\sum_{i\in I_W}a_iz_i=:u+w.
\]
$U,W$는 부분공간이므로 $u\in U$, $w\in W$. 따라서 $y\in U+W$. 즉
\[
\operatorname{span}(U\cup W)\subset U+W.
\]
두 포함으로 등식이 성립한다.
\end{proof}

\paragraph{의미 해석.}
이 정리는 합공간을 ``두 부분공간에서 온 벡터들의 선형결합 전체''로 해석하게 해 줍니다. 따라서 합공간 문제를 생성공간 문제로 바꾸어 다룰 수 있고, 계산 전략이 크게 단순해집니다.

\begin{example}
$P_3(\R)$에서
\[
E=\{a+ct^2:a,c\in\R\},\qquad O=\{bt+dt^3:b,d\in\R\}
\]
로 두자. 임의의
\[
p(t)=a+bt+ct^2+dt^3
\]
에 대해
\[
p(t)=(a+ct^2)+(bt+dt^3)\in E+O
\]
이므로 $P_3(\R)=E+O$. 또한 $E\cap O=\{0\}$이므로 $P_3(\R)=E\oplus O$이다.
\end{example}

\subsection*{주의}
\begin{warning}
$U+W$는 집합의 합집합 $U\cup W$와 다르다. $U+W$에는 $u+w$ 형태의 새로운 벡터가 포함된다.
\end{warning}
\begin{warning}
$U\cap W=\{0\}$만으로는 직합이 완성되지 않는다. 반드시 $V=U+W$도 함께 확인해야 한다.
\end{warning}

\subsection*{자가진단퀴즈}
\begin{enumerate}[label=\arabic*.]
\item $\R^2$에서 $U=\operatorname{span}\{(1,0)\}$, $W=\operatorname{span}\{(1,1)\}$일 때 $\R^2=U\oplus W$임을 보이라.
\item $\R^3$에서 $U=\operatorname{span}\{(1,0,0),(0,1,0)\}$, $W=\operatorname{span}\{(1,1,0)\}$일 때 $U+W$와 $U\cap W$를 구하라.
\item $V=U\oplus W$이면 $U$와 $W$의 임의 원소 $u,w$에 대해 $u+w=0$이면 $u=w=0$임을 증명하라.
\end{enumerate}

\chapterapplicationtemplate{
신호 분해를 직합의 관점에서 정리해 보겠습니다. $V=C([-1,1],\R)$라 두고
\[
E=\{f\in V:f(-x)=f(x)\},\qquad
O=\{f\in V:f(-x)=-f(x)\}
\]
를 정의합니다. 임의의 $f\in V$에 대해
\[
f_e(x)=\frac{f(x)+f(-x)}{2},\qquad
f_o(x)=\frac{f(x)-f(-x)}{2}
\]
로 두면 $f_e\in E$, $f_o\in O$, 그리고 $f=f_e+f_o$가 성립합니다. 또한 $E\cap O=\{0\}$이므로 $V=E\oplus O$입니다. 따라서 임의의 신호를 짝성분과 홀성분으로 유일하게 분해할 수 있고, 대칭성 분석이나 필터 설계에서 이를 직접 사용할 수 있습니다.
}

\chaptersummarytemplate

\chapterexercisestemplate

\subsection*{연습문제 A (기초 확인)}
\begin{enumerate}[label=A\arabic*.]
\item $\R^3$의 부분집합
\[
W=\{(x,y,z)\in\R^3:x-2y+z=0\}
\]
가 부분공간인지 판정하라.
\item $\R^2$의 부분집합
\[
S=\{(x,y)\in\R^2:x+y=1\}
\]
가 부분공간이 아닌 이유를 한 줄 공리 점검으로 설명하라.
\item $V$가 벡터공간일 때 임의의 $v\in V$에 대해 $0_Kv=0_V$를 다시 증명하라.
\item $S=\{(1,0,1),(0,1,1)\}\subset\R^3$에 대해 $\operatorname{span}(S)$를 매개변수 형태로 쓰라.
\item 벡터열 $(1,2),(2,4),(3,0)$의 일차독립 여부를 판정하라.
\item $U=\operatorname{span}\{(1,0,0)\}$, $W=\operatorname{span}\{(0,1,0)\}$일 때 $U+W$를 구하라.
\item $P_2(\R)$에서 $\{1,t,t^2\}$가 생성집합임을 보이라.
\item $\R^2$에서 $U=\operatorname{span}\{(1,0)\}$, $W=\operatorname{span}\{(0,1)\}$가 직합을 이루는지 판정하라.
\end{enumerate}

\subsection*{연습문제 B (개념 연결)}
\begin{enumerate}[label=B\arabic*.]
\item $A\in M_{m,n}(K)$에 대하여 해집합 $\{x\in K^n:Ax=0\}$가 부분공간임을 부분공간 판정법으로 증명하라.
\item 임의의 부분집합 $S,T\subset V$에 대해
\[
\operatorname{span}(S\cup T)=\operatorname{span}(S)+\operatorname{span}(T)
\]
를 증명하라.
\item 유한 집합 $S\subset V$가 일차종속이면 어떤 $v\in S$가 존재하여
\[
\operatorname{span}(S)=\operatorname{span}(S\setminus\{v\})
\]
임을 증명하라.
\item $\R^3$에서
\[
U=\operatorname{span}\{(1,0,0),(0,1,0)\},\quad
W=\operatorname{span}\{(1,1,\lambda)\}
\]
에 대해 $\R^3=U\oplus W$가 되기 위한 $\lambda$ 조건을 구하라.
\item $P_3(\R)$에서 짝다항식 부분공간 $E$, 홀다항식 부분공간 $O$를 정의하고 $P_3(\R)=E\oplus O$를 증명하라.
\end{enumerate}

\subsection*{연습문제 C (증명/종합)}
\begin{enumerate}[label=C\arabic*.]
\item $V=U\oplus W$라 하자. 각 $v\in V$의 유일분해 $v=u+w$를 이용해 $P_U(v)=u$를 정의할 때, $P_U$가 선형사상이고
\[
P_U^2=P_U,\qquad \ker P_U=W,\qquad \operatorname{im}P_U=U
\]
임을 증명하라.
\item 유한 벡터열 $(v_1,\dots,v_n)$이 일차독립일 필요충분조건이
\[
v_j\notin\operatorname{span}(\{v_i:i\neq j\})\quad(1\le j\le n)
\]
임을 증명하라.
\item 부분공간 $U,W,Z\le V$에 대해 $V=U\oplus W$이고 $Z\subset U$이면
\[
Z\cap W=\{0_V\},\qquad Z\oplus W\subset U\oplus W
\]
를 증명하라.
\end{enumerate}